\section{Contoh Penerapan}
\label{sec:contoh-bab2}

\subsection{Contoh 1: Konversi Bilangan Bertahap}

Berikut contoh konversi yang dikerjakan tanpa alat bantu, dengan
langkah yang dapat diperiksa.

\begin{enumerate}
    \item \textbf{Desimal 250 ke biner.}
          Bagi berulang oleh 2 dan catat sisanya:
          $250 = 1111\,1010_2$.

    \item \textbf{Biner \bin{1111\,1010} ke heksadesimal.}
          Kelompokkan per empat bit dari kanan: \bin{1111} dan
          \bin{1010}, yaitu \hex{F} dan \hex{A}. Hasilnya
          \hex{FA}.

    \item \textbf{Heksadesimal \hex{FA} ke desimal.}
          $15 \times 16 + 10 = 240 + 10 = 250$. Hasilnya konsisten
          dengan langkah pertama, sehingga ketiga konversi saling
          memverifikasi.
\end{enumerate}

\subsection{Contoh 2: Menghitung \eng{Two's Complement}}

Untuk menyatakan \dec{-25} dalam 8 bit:

\begin{enumerate}
    \item Tulis nilai positifnya dalam 8 bit: \bin{0001\,1001}
          (yaitu \hex{19}).
    \item Balikkan seluruh bitnya: \bin{1110\,0110} (yaitu
          \hex{E6}).
    \item Tambahkan satu: \bin{1110\,0111} (yaitu \hex{E7}).
\end{enumerate}

Verifikasi: \hex{E7} + \hex{19} = \hex{100}, yang dalam 8 bit menjadi
\hex{00} dengan \eng{carry} keluar. Hasilnya nol, sebagaimana
mestinya untuk $(-25) + 25$. Sifat inilah yang membuat
\eng{two's complement} dipakai: prosesor tidak memerlukan rangkaian
khusus untuk pengurangan, cukup penjumlahan biasa.

\subsection{Contoh 3: Melihat Penyandian Data di Memori}

Program pada Kode~\ref{lst:penyandian} mendeklarasikan beberapa data
dengan penyandian berbeda, lalu mahasiswa memeriksa bagaimana
byte-byte tersebut benar-benar tersimpan.

\begin{lstlisting}[
    language=MASM,
    caption={Program yang memperlihatkan penyandian data di memori.},
    label={lst:penyandian}
]
.MODEL SMALL
.STACK 100h

.DATA
    nilai_desimal   DB  25          ; nilai biner, bukan sandian karakter
    nilai_bcd       DB  25h         ; 25 dalam BCD
    sandian_karakter DB '2'         ; sandian ASCII untuk karakter '2'
    kata_16bit      DW  1234h       ; dua byte, little-endian

.CODE
START:
    MOV AX, @DATA
    MOV DS, AX

    MOV AH, 4Ch
    MOV AL, 00h
    INT 21h

END START
\end{lstlisting}

Setelah program dirakit, periksa isi segmen data dengan debugger.
Tabel~\ref{tab:penyandian-hasil} memperlihatkan hasil yang
seharusnya Anda temukan.

\begin{table}[H]
    \centering
    \caption{Isi memori untuk data pada Kode~\ref{lst:penyandian}.}
    \label{tab:penyandian-hasil}
    \small
    \begin{tabular}{@{}llp{7.2cm}@{}}
        \toprule
        \textbf{Data} & \textbf{Isi memori} & \textbf{Penjelasan} \\
        \midrule
        \code{nilai\_desimal} & \hex{19} &
        \dec{25} dalam biner; jika dicetak langsung, yang muncul
        adalah karakter tak terlihat, bukan ``25'' \\
        \addlinespace
        \code{nilai\_bcd} & \hex{25} &
        Dua digit BCD dalam satu byte: \bin{0010} dan \bin{0101} \\
        \addlinespace
        \code{sandian\_karakter} & \hex{32} &
        Sandian ASCII karakter \code{'2'}, yaitu \hex{30} + 2 \\
        \addlinespace
        \code{kata\_16bit} & \hex{34} lalu \hex{12} &
        Nilai \hex{1234} disimpan \eng{little-endian}: byte rendah
        lebih dahulu \\
        \bottomrule
    \end{tabular}
\end{table}

\begin{infobox}
Perhatikan tiga baris pertama: ketiganya ``bernilai 25'' menurut
manusia, tetapi tiga-tiganya tersimpan berbeda di memori. Inilah
sebabnya mahasiswa pemrograman tingkat rendah harus selalu
memikirkan \emph{representasi}, bukan hanya \emph{nilai}.
\end{infobox}

\subsection{Contoh 4: Menjumlahkan Data dan Memeriksa Flag}

Program pada Kode~\ref{lst:flag-uji} memperlihatkan bagaimana
representasi yang sama menghasilkan tafsiran berbeda.

\begin{lstlisting}[
    language=MASM,
    caption={Penjumlahan yang menghasilkan \eng{carry} dan
    \eng{overflow} sekaligus.},
    label={lst:flag-uji}
]
.MODEL SMALL
.STACK 100h

.CODE
START:
    MOV AX, @DATA
    MOV DS, AX

    MOV AL, 0FFh        ; 255 tak bertanda, atau -1 bertanda
    ADD AL, 01h         ; 255 + 1

    MOV AH, 4Ch
    MOV AL, 00h
    INT 21h

END START
\end{lstlisting}

Setelah \asm{ADD AL, 01h} dieksekusi: \reg{AL} bernilai \hex{00},
\flag{CF} bernilai 1 (ada \eng{carry} keluar dari bit ke-8),
\flag{ZF} bernilai 1 (hasilnya nol), dan \flag{OF} bernilai 0
(tidak ada \eng{overflow} bertanda karena $-1 + 1 = 0$ masih dalam
rentang). Tafsiran \eng{carry} dan \eng{overflow} bergantung pada
apakah program bermaksud memperlakukan \hex{FF} sebagai
\dec{255} atau sebagai $-1$.

\begin{session}[caption={Pengamatan register dan flag pada
Kode~\ref{lst:flag-uji}.}]
-R
AX=00FF  BX=0000  CX=0000  DX=0000  SP=FFFE  BP=0000  SI=0000  DI=0000
DS=0700  ES=0700  SS=0700  CS=0700  IP=0106   NV UP EI PL NZ NA PO NC
-T
AX=0000  BX=0000  CX=0000  DX=0000  SP=FFFE  BP=0000  SI=0000  DI=0000
DS=0700  ES=0700  SS=0700  CS=0700  IP=0108   NV UP EI PL ZR AC PE CY
\end{session}

Baris pertama memperlihatkan keadaan sebelum \asm{ADD AL, 01h}:
\reg{AL} berisi \hex{FF} dan seluruh flag utama bernilai 0, yang pada
DEBUG ditandai \code{NV} (tidak \eng{overflow}), \code{NZ} (tidak
nol), dan \code{NC} (tidak ada \eng{carry}).

Baris kedua memperlihatkan keadaan sesudahnya. Penanda berubah
menjadi \code{ZR} (\flag{ZF} = 1, hasilnya nol), \code{CY}
(\flag{CF} = 1, ada \eng{carry} keluar bit ke-8), dan \code{AC}
(\flag{AF} = 1, ada \eng{carry} dari \eng{nibble} rendah ke tinggi).
Penanda \code{NV} tetap, yang berarti \flag{OF} = 0: sebagai
bilangan bertanda, $-1 + 1 = 0$ tidak melampaui rentang sehingga
tidak terjadi \eng{overflow}.

Inilah contoh paling ringkas mengapa \eng{carry} dan
\eng{overflow} harus dibedakan: pada operasi yang sama,
\flag{CF} bernilai 1 sementara \flag{OF} bernilai 0.
